\section*{第 \ref{ch:b2:curves} 章 --- 曲线}

\begin{solution}{exo:b2:curves:1}
$\gamma'(t) = (1 - \cos t,\ \sin t)$，于是
\[
\norm{\gamma'(t)}^2 = (1 - \cos t)^2 + \sin^2 t
= 2 - 2\cos t = 4\sin^2\tfrac t2 ,
\]
且 $\norm{\gamma'(t)} = 2\sin\frac t2$（在 $[0,
2\pi]$ 上非负）。由此可得
\[
L = \int_0^{2\pi} 2\sin\tfrac t2\,\dd t
= \Bigl[-4\cos\tfrac t2\Bigr]_0^{2\pi} = 8 :
\]
摆线的一拱长为 $8$（轮子半径为
$1$）--- 这是雷恩的一个著名结果，其中不出现 $\pi$。
\end{solution}

\begin{solution}{exo:b2:curves:2}
设 $x = a\cos t$，$y = b\sin t$，则 $x' = -a\sin t$，$y' = b\cos
t$，$x'' = -a\cos t$，$y'' = -b\sin t$，由
\cref{prop:b2:curves:kappaformula}可得
\[
\kappa(t)
= \frac{x'y'' - y'x''}{(x'^2 + y'^2)^{3/2}}
= \frac{ab\sin^2 t + ab\cos^2 t}
       {(a^2\sin^2 t + b^2\cos^2 t)^{3/2}}
= \frac{ab}{(a^2\sin^2 t + b^2\cos^2 t)^{3/2}} .
\]
分母在 $\sin t = 0$ 时取最小值（值为 $b^3$，点为
$(\pm a, 0)$），在 $\cos t = 0$ 时取最大值（值为 $a^3$，点为
$(0, \pm b)$）。由 $a > b$，可知 $\kappa$ 在
长轴端点处取最大值 $\kappa_{\max} = a/b^2$，并在
短轴端点处取最小值 $\kappa_{\min} = b/a^2$：椭圆
在长轴端点处弯曲得最厉害。
\end{solution}

\begin{solution}{exo:b2:curves:3}
对图像 $\gamma(x) = (x, \cosh x)$，有 $\norm{\gamma'(x)} =
\sqrt{1 + \sinh^2 x} = \cosh x$，于是从 $0$ 到 $x$ 的弧长
为 $\int_0^x \cosh u\,\dd u = \sinh x$。图像的曲率
由\cref{prop:b2:curves:kappaformula}给出：
\[
\kappa(x) = \frac{f''(x)}{(1 + f'(x)^2)^{3/2}}
= \frac{\cosh x}{\cosh^3 x} = \frac{1}{\cosh^2 x} ,
\]
故 $R(x) = \cosh^2 x$，与所述一致。值得注意的是
$R(x) = 1 + s(x)^2$，其中 $s = \sinh x$ 为弧长：
悬链线的曲率半径随弧长的平方
增长，弧长自顶点起算。
\end{solution}

\begin{solution}{exo:b2:curves:4}
$\gamma'(t) = e^t(\cos t - \sin t,\ \sin t + \cos t)$，于是
\[
\langle \gamma(t), \gamma'(t)\rangle = e^{2t}
\bigl(\cos t(\cos t - \sin t) + \sin t(\sin t + \cos t)\bigr)
= e^{2t},
\]
而 $\norm{\gamma(t)} = e^t$，$\norm{\gamma'(t)} =
e^t\sqrt 2$。于是
\[
\cos\angle\bigl(\gamma, \gamma'\bigr)
= \frac{e^{2t}}{e^t \cdot e^t\sqrt2} = \frac{1}{\sqrt2} :
\]
切线与向径始终成 $\pi/4$ 的夹角 --- 这是对数螺线的
\emph{等角}性质。弧长
在 $(-\infty, 0]$ 上为：
\[
\int_{-\infty}^0 \norm{\gamma'(t)}\,\dd t
= \sqrt2\int_{-\infty}^0 e^t\,\dd t = \sqrt 2 ,
\]
弧长有限，尽管螺线绕
原点旋转了无穷多圈。曲率方面，$x'y'' - y'x''$ 由 $\gamma'' =
e^t(-2\sin t,\ 2\cos t)$，得
\[
x'y'' - y'x'' = e^{2t}\bigl(2\cos t(\cos t - \sin t)
+ 2\sin t(\sin t + \cos t)\bigr) = 2e^{2t},
\]
于是 $\kappa(t) = \dfrac{2e^{2t}}{(e^t\sqrt2)^3} =
\dfrac{1}{e^t\sqrt2}$。曲率为 $1/(\sqrt2\,
\norm{\gamma})$，随螺线增大而衰减。
\end{solution}

\begin{solution}{exo:b2:curves:5}
\emph{第一支：}$\gamma(t) = (t^2,\ t^4 + t^5)$。
在 $0$ 处，$\gamma'' = (2, 0) \neq 0$，故 $p = 2$。接着 $\gamma^{(3)}(0)
= (0, 0)$，$\gamma^{(4)}(0) = (0, 24)$，与 $(2,
0)$ 不共线，故 $q = 4$。两者均为偶数：\emph{第二类尖点} --- 两支
均沿 $+u = (1,0)$ 方向离开，
并位于切线的同一侧。（事实上，在两支上 $y = x^2 \pm x^{5/2}$，
当 $x$ 充分小时符号相同。）

\emph{第二支：}$\gamma(t) = (t^3, t^4)$。$\gamma'(0) =
\gamma''(0) = 0$，$\gamma^{(3)}(0) = (6, 0)$，故 $p = 3$，为奇数。下一步
$\gamma^{(4)}(0) = (0, 24)$，故 $q = 4$，为偶数。一奇一偶时
为\emph{普通点} --- 尽管速度为零，
轨迹 $y = x^{4/3}$ 平滑地穿过原点，
并保持在其切线 $y = 0$ 的上方。
\end{solution}

\begin{solution}{exo:b2:curves:6}
对 $c(s) = \gamma(s) + \dfrac{1}{\kappa(s)}N(s)$ 求导，利用
平面弗雷内公式
（\cref{thm:b2:curves:frenet2d}）：
\[
c'(s) = T(s) - \frac{\kappa'(s)}{\kappa(s)^2}N(s)
+ \frac{1}{\kappa(s)}\,\bigl(-\kappa(s)T(s)\bigr)
= -\frac{\kappa'(s)}{\kappa(s)^2}\,N(s) ,
\]
切向项恰好相消。于是渐屈线的速度
由 $N(s)$ 承载，它指向
$\gamma$ 在 $\gamma(s)$ 处的法线 --- 而点 $c(s)$ 恰位于
该法线上：渐屈线即为\emph{法线的包络}。
（在 $\kappa' = 0$ 处渐屈线有奇点，这正是
椭圆渐屈线出现尖点的原因。）
\end{solution}

\begin{solution}{exo:b2:curves:7}
设 $G(s) = F(s)F(s)^{\mathsf T}$，利用 $F' = FA$ 及
$(F^{\mathsf T})' = (F')^{\mathsf T} = A^{\mathsf T}F^{\mathsf T}$,
\[
G' = F'F^{\mathsf T} + F(F^{\mathsf T})'
= FAF^{\mathsf T} + FA^{\mathsf T}F^{\mathsf T}
= F(A + A^{\mathsf T})F^{\mathsf T} = 0
\]
由反对称性可知 $G$ 在区间上为常值，等于
$G(s_0) = F(s_0)F(s_0)^{\mathsf T} = I$：$F(s)$ 对每个 $s$ 都是正交矩阵。
（另证：不直接算 $G'$，而注意 $G$ 与常值 $I$ 两者
都是线性方程组 $Y' = YA +
A^{\mathsf T}Y$ 的解，初值相同。由线性方程组的柯西–利普希茨唯一性
（见\cref{ch:b2:diffeq}）可得 $G
\equiv I$。）

\emph{关联：}弗雷内方程组 $(T, N, B)' = (T, N, B)\,A(s)$
的系数矩阵为如下反对称矩阵
\[
A = \begin{pmatrix} 0 & -\kappa & 0\\ \kappa & 0 & -\tau\\
0 & \tau & 0\end{pmatrix}
\]
（矩阵的各列分别给出 $T', N', B'$ 的表示）。上述计算表明
即：从标准正交出发的解标架 \emph{始终}
保持标准正交 --- 这正是由 $(\kappa, \tau)$ 重建曲线的
基本定理中的关键一步。
\end{solution}

\begin{solution}{exo:b2:curves:8}
由\cref{thm:b2:curves:fundamental}（唯一性部分）可知，存在一个
$\mathcal{C}^1$ 类角度函数 $\varphi$，使得 $T(s) =
(\cos\varphi(s), \sin\varphi(s))$ 且 $\varphi' = \kappa$。于是
\[
\int_0^L \kappa(s)\,\dd s = \varphi(L) - \varphi(0) .
\]
由于该弧是周期为 $L$ 的闭弧，故 $T(L) = T(0)$：
$(\cos\varphi(L), \sin\varphi(L)) = (\cos\varphi(0),
\sin\varphi(0))$，故 $\varphi(L) - \varphi(0) \in 2\pi\Z$。闭曲线的
全曲率因而总是一个整数
$2\pi$ 的整数倍 --- 该整数即 \emph{旋转数}
（即切向量 $T$ 所转过的整圈数）。对于
半径为 $R$ 的圆，$\kappa = 1/R$，$L = 2\pi R$，全曲率
为 $2\pi$，旋转数为 $1$；转向切线
定理断言该值对每条简单闭曲线都成立。
\end{solution}

\begin{solution}{exo:b2:curves:9}
由于 $\tau \equiv 0$，曲线位于某一平面内
（见\cref{prop:b2:curves:torsion}），在该平面内讨论。考虑
候选圆心
\[
c(s) = \gamma(s) + \frac{1}{\kappa}N(s)
\qquad (\kappa \text{ 为常数}).
\]
利用弗雷内公式求导（$N' = -\kappa T + \tau B
= -\kappa T$，此时）：
\[
c'(s) = T + \frac1\kappa(-\kappa T) = 0 ,
\]
故 $c$ 为定点 $\Omega$。则 $\norm{\gamma(s) -
\Omega} = \norm{-\frac1\kappa N(s)} = \frac1\kappa$，对一切 $s$：
曲线位于圆心为 $\Omega$、半径为
$1/\kappa$ 的圆上（位于其所在平面内），
且它是该圆的非常值弧，故为该圆的一段弧。
\end{solution}

\begin{solution}{exo:b2:curves:10}
对于图像 $\gamma(x) = (x, x^2/2)$，$\norm{\gamma'(x)} =
\sqrt{1 + x^2}$，故 $L = \int_0^a\sqrt{1+x^2}\,\dd x$。
作代换 $x = \sinh u$（$\dd x = \cosh u\,\dd u$，$u$ 从
$0$ 变到 $u_a = \ln(a + \sqrt{1+a^2})$）：
\[
L = \int_0^{u_a}\cosh^2 u\,\dd u
= \frac12\bigl[u + \sinh u\cosh u\bigr]_0^{u_a}
= \frac12\Bigl(\ln\bigl(a + \sqrt{1+a^2}\bigr) +
a\sqrt{1+a^2}\Bigr),
\]
利用 $\cosh^2 u = \frac{1 + \cosh 2u}2$ 及 $\sinh u_a = a$，
$\cosh u_a = \sqrt{1 + a^2}$.
\end{solution}

\begin{solution}{exo:b2:curves:11}
按弧长参数化
（见\cref{thm:b2:curves:arclength}），并令 $\lambda(s) =
\langle P - \gamma(s), T(s)\rangle$，它是一个 $\mathcal C^1$
函数。由于 $P$ 位于 $\gamma(s)$ 处的切线上，
向量 $P - \gamma(s)$ 与 $T(s)$ 共线，故 $P =
\gamma(s) + \lambda(s)T(s)$。求导得
\[
0 = T(s) + \lambda'(s)T(s) + \lambda(s)T'(s)
= \bigl(1 + \lambda'(s)\bigr)T(s) + \lambda(s)T'(s),
\]
又 $T'(s) \perp T(s)$（对 $\norm T^2 = 1$ 求导），故
两个分量均为零：$\lambda' = -1$ 且 $\lambda T' = 0$。
于是 $\lambda(s) = c - s$ 至多一次为零，因而 $T' = 0$
在稠密集上成立，由连续性知处处成立：$T$ 为
常值单位向量，且 $\gamma(s) = \gamma(s_0) + (s -
s_0)T$：一条直线（必过 $P$）。
\end{solution}

\begin{solution}{exo:b2:curves:12}
(a) 弗雷内矩阵
\[
A(s) = \begin{pmatrix} 0 & -\kappa & 0\\ \kappa & 0 & -\tau\\
0 & \tau & 0\end{pmatrix}
\]
其元素连续，故线性方程组 $F' = FA(s)$，
$F(s_0) = F_0$（正向标准正交矩阵）的初值问题有唯一
解，且该解在整个 $J$ 上存在（见\cref{thm:b2:diffeq:linear}）。
由\cref{exo:b2:curves:7}可知 $F(s)$ 对每个 $s$ 都是正交矩阵；
$\det F$ 连续，取值于 $\{\pm1\}$，且等于 $1$
（在 $s_0$ 处），故 $F(s)$ 对一切 $s$ 为正向矩阵。

(b) 读出 $T, N, B$ 为 $F$ 的各行（使得 $T' = \kappa
N$，$N' = -\kappa T + \tau B$，$B' = -\tau N$）并令
$\gamma(s) = \gamma_0 + \int_{s_0}^s T(u)\,\dd u$。则
$\gamma' = T$ 为单位向量，即单位速率；$T' = \kappa N$
其中 $\kappa > 0$，且 $N$ 为与 $T$ 正交的单位向量，故 $\gamma$
是双正则曲线，曲率为 $\norm{T'} = \kappa$，
主法向量为 $N$；副法向量为 $T \wedge N = B$
（正向标准正交标架），且 $B' = -\tau N$ 确定
挠率为 $\tau$。

(c) 设 $\gamma_1, \gamma_2$ 为单位速率双正则曲线，
且具有相同的 $(\kappa, \tau)$。存在唯一的正向
等距 $\Phi = \rho + w$（$\rho \in SO(3)$），将
$\gamma_1(s_0)$ 映为 $\gamma_2(s_0)$，并将
$\gamma_1$ 在 $s_0$ 处的弗雷内标架映为 $\gamma_2$ 在 $s_0$ 处的标架。
曲线 $\Phi\circ\gamma_1$ 为单位速率曲线，具有相同的不变量
（其标架为将 $\rho$ 作用于 $\gamma_1$ 的标架所得，
$\rho$ 保持叉积，因其为正向）。现在
$\Phi\circ\gamma_1$ 与 $\gamma_2$ 的标架都满足 $F' =
FA(s)$ 初值相同，由唯一性知二者重合；
特别地切向量一致，从公共点 $s_0$ 积分
得 $\Phi\circ\gamma_1 = \gamma_2$。
\end{solution}

\begin{solution}{pb:b2:curves:1}
\textbf{1.} 该方程组关于 $(x, y)$ 是线性的，行列式
$\Delta(t) = a b' - a'b \neq 0$，由克莱姆法则得
唯一解：
\[
x = \frac{c b' - c' b}{\Delta}, \qquad
y = \frac{a c' - a' c}{\Delta}.
\]

\textbf{2.} 由于 $a(t)x(t) + b(t)y(t) = c(t)$ 恒成立，
求导得 $a'x + b'y + ax' + by' = c'$，由
第二个特征方程使 $a'x + b'y - c'$ 为零，故
$a(t)x'(t) + b(t)y'(t) = 0$，即 $E'(t)$ 与 $(a,
b)$ 正交，从而与 $(-b, a)$ 平行，即 $D_t$ 的方向。由于
$E(t) \in D_t$（第一个方程），且 $E'(t) \neq 0$，直线
$D_t$ 恰为曲线 $E$ 在 $E(t)$ 处的切线。

\textbf{3.} 此时 $(a, b, c) = (t, -1, t^2/2)$，故 $\Delta =
t\cdot0 - 1\cdot(-1) = 1$，且
\[
x = \frac{c b' - c' b}{\Delta} = \tfrac{t^2}2\cdot 0 +
t = t, \qquad
y = \frac{a c' - a' c}{\Delta} = t\cdot t - \tfrac{t^2}2 =
\tfrac{t^2}2 :
\]
抛物线切线族的包络即
抛物线本身，与预期相符。

\textbf{4.} 法线为 $\langle M, T(s)\rangle =
\langle\gamma(s), T(s)\rangle$，系数为 $a = T_1$，$b =
T_2$，$c = \langle\gamma, T\rangle$。求导并利用
弗雷内公式（$T' = \kappa N$）得 $a' = \kappa N_1$，$b' = \kappa
N_2$，且 $c' = \langle T, T\rangle + \langle\gamma, \kappa
N\rangle = 1 + \kappa\langle\gamma, N\rangle$。第二个
特征方程 $\kappa\langle M, N\rangle = 1 +
\kappa\langle\gamma, N\rangle$ 可写为 $\kappa\langle M -
\gamma, N\rangle = 1$。第一式表明 $M - \gamma \perp T$，
故 $M - \gamma = \mu N$，$\mu = 1/\kappa$，于是特征点
为 $\gamma + \frac1\kappa N$，即曲率中心，
而法线族的包络即渐屈线（见
\cref{exo:b2:curves:6}）。最后 $\Delta = T_1\kappa N_2 -
\kappa N_1 T_2 = \kappa\det(T, N) = \kappa \neq 0$.

\textbf{5.} 线束：方程组 $x\cos\theta + y\sin\theta =
0$，$-x\sin\theta + y\cos\theta = 0$ 的行列式为 $1$，且
解恒为 $(0,0)$（对每个 $\theta$）：$E \equiv (0,0)$，$E'
\equiv 0$，故不存在曲线 --- 只是所有直线的公共点。
平行直线族：$a' = b' = 0$ 给出
$\Delta \equiv 0$，第二个方程为 $0 = c'(t)$，
一旦直线族真正运动便不成立：没有特征
点，事实上平行直线族沿其全部成员都不与
任何曲线相切。

\textbf{6.} 过 $(\cos t, 0)$ 与 $(0, \sin t)$ 的直线
为 $\frac x{\cos t} + \frac y{\sin t} = 1$，即 $x\sin t +
y\cos t = \sin t\cos t$。此时 $(a, b, c) = (\sin t, \cos t,
\sin t\cos t)$：$a' = \cos t$，$b' = -\sin t$，$c' = \cos
2t$，$\Delta = -\sin^2 t - \cos^2 t = -1$。由克莱姆法则：
\begin{align*}
x &= \frac{cb' - c'b}{-1} = \sin^2 t\cos t + \cos 2t\cos t
= \cos t\,(\sin^2 t + \cos^2 t - \sin^2 t) = \cos^3 t,\\
y &= \frac{ac' - a'c}{-1} = \sin t\cos^2 t - \sin t\cos 2t
= \sin t\,(\cos^2 t - \cos^2 t + \sin^2 t) = \sin^3 t .
\end{align*}
又 $(\cos^3t)^{2/3} + (\sin^3t)^{2/3} = 1$，此即星形线。

\textbf{7.} $E'(t) = 3(-\cos^2 t\sin t,\ \sin^2 t\cos t)$
在 $t = 0$ 处为零。此处 $E''(0) = (-3, 0) \neq 0$，
于是 $p = 2$；$x$ 分量（取自 $E$）关于 $t$ 为偶，故
$E'''(0) = (0, 6)$，与前一阶导数不共线，$q = 3$。由奇偶性，
$(1, 0)$ 处为第一类尖点
（见\cref{prop:b2:curves:local}），切线沿
$x$ 轴方向。星形线的对称变换 $x \mapsto -x$，$y \mapsto -y$，$(x,
y) \mapsto (y, x)$ 将尖点分别移至 $(-1,
0)$ 与 $(0, \pm1)$。

\textbf{8.} $E'(t) = 3\sin t\cos t\,(-\cos t, \sin t)$，故
$\norm{E'(t)} = 3\abs{\sin t\cos t} = \tfrac32\abs{\sin 2t}$.
在一个象限上：$\int_0^{\pi/2}\tfrac32\sin 2t\,\dd t =
\tfrac32$，由对称性，总长为 $4 \cdot
\tfrac32 = 6$.

\textbf{9.} $E(t) - P_t = (\cos^3 t - \cos t,\ \sin^3 t) =
\sin^2 t\,(-\cos t,\ \sin t) = \sin^2 t\,(Q_t - P_t)$。因此
接触点为 $(P_t, \cos^2 t)$ 与
$(Q_t, \sin^2 t)$ 的重心：当 $t$ 由 $0$ 增至 $\pi/2$ 时，
它从梯子的地面端滑向墙端。

\textbf{10.} 在第一象限中，星形线下方的区域
面积为 $\int_0^1 y\,\dd x$，其中 $x = \cos^3 t$
由 $1$ 减至 $0$，$t$ 从 $0$ 增至 $\pi/2$：
\[
\int_0^1 y\,\dd x
= \int_{\pi/2}^{0}\sin^3 t\,(-3\cos^2 t\sin t)\,\dd t
= 3\int_0^{\pi/2}\sin^4 t\cos^2 t\,\dd t .
\]
降幂化简：$\sin^4 t\cos^2 t = (\sin t\cos t)^2\sin^2 t =
\tfrac18\bigl(\sin^2 2t - \sin^2 2t\cos 2t\bigr)$，且
$\int_0^{\pi/2}\sin^2 2t\,\dd t = \tfrac\pi4$，而
$\int_0^{\pi/2}\sin^2 2t\cos 2t\,\dd t =
\bigl[\tfrac{\sin^3 2t}6\bigr]_0^{\pi/2} = 0$。于是
积分为 $\tfrac\pi{32}$，第一象限面积为
$\tfrac{3\pi}{32}$，封闭区域面积为 $4 \cdot
\tfrac{3\pi}{32} = \tfrac{3\pi}8$。

\textbf{11.} 在 $\gamma(t) = (t, t^2/2)$ 处，切线
方向为 $(1, t)$，故法线为 $\{(x, y) : (x -
t) + t\,(y - t^2/2) = 0\}$，即 $x + t\,y = t +
\tfrac{t^3}2$。

\textbf{12.} $(a, b, c) = (1, t, t + t^3/2)$：$a' = 0$，$b' =
1$，$c' = 1 + \tfrac32 t^2$，$\Delta = 1$。则 $y = ac' -
a'c = 1 + \tfrac32t^2$，且
\[
x = cb' - c'b = t + \tfrac{t^3}2 - t\Bigl(1 +
\tfrac32t^2\Bigr) = -t^3 .
\]
消去 $t$：$t^2 = \tfrac23(y - 1)$，且 $x^2 = t^6 =
\tfrac8{27}(y - 1)^3$：一条半立方抛物线，顶点为
$(0, 1)$。

\textbf{13.} $T = (1, t)/\sqrt{1+t^2}$，$N = (-t,
1)/\sqrt{1+t^2}$，且 $\kappa = (1+t^2)^{-3/2}$，于是
\[
\gamma + \frac1\kappa N = (t, \tfrac{t^2}2) +
(1+t^2)\,(-t, 1) = \Bigl(-t^3,\ 1 + \tfrac32t^2\Bigr),
\]
与原曲线相同：法线的包络是
曲率中心的轨迹，正如第 4 题所言。

\textbf{14.} $E'(t) = (-3t^2, 3t)$ 在 $t = 0$ 处为零；
$E''(0) = (0, 3) \neq 0$ 得 $p = 2$，且 $E'''(0) = (-6,
0)$ 得 $q = 3$：$(0, 1)$ 处为第一类尖点。
顶点处 $\kappa = (1+t^2)^{-3/2}$ 取最大值，故
$\kappa'(0) = 0$，且渐屈线的速度
$-\frac{\kappa'}{\kappa^2}N$ 恰在该处为零，渐屈线的尖点
位于曲率的极值点处。

\textbf{15.} 参数为 $t$ 的法线经过
$(x_0, y_0)$ 当且仅当 $x_0 + t\,y_0 = t + \tfrac{t^3}2$，即
\[
\frac{t^3}2 + (1 - y_0)\,t - x_0 = 0 ,
\]
关于 $t$ 的三次方程，有一个或三个实根（不计重数，
对一般点而言）。在轴 $x_0 = 0$ 上，它分解为
$t\bigl(\tfrac{t^2}2 + 1 - y_0\bigr) = 0$，其根为
$t = 0$（轴即为顶点处的法线），以及
$t = \pm\sqrt{2(y_0 - 1)}$，当 $y_0 > 1$ 时。综上，
$y_0 < 1$ 时只有一条法线，$y_0 > 1$ 时有三条，而在 $y_0 = 1$ 处的
三重根标志着渐屈线的尖点。一般地，
三次方程的二重根意味着该点同时满足
直线方程及其关于 $t$ 的导数 --- 即该点\emph{位于
包络上}：渐屈线恰为
一条法线区域与三条法线区域之间的分界。

\textbf{16.} 镜面反射使入射方向的法向分量反号，
切向分量保持不变：
记 $u = \langle u, n\rangle n + u_{\mathrm{tan}}$，
则反射方向为 $u_{\mathrm{tan}} - \langle u,
n\rangle n = u - 2\langle u, n\rangle n$。此处 $\langle u,
n\rangle = \cos\theta$，故
\[
v = (1, 0) - 2\cos\theta\,(\cos\theta, \sin\theta)
= (1 - 2\cos^2\theta,\ -2\sin\theta\cos\theta)
= -(\cos2\theta,\ \sin2\theta).
\]

\textbf{17.} 反射光线经过 $P_\theta =
(\cos\theta, \sin\theta)$，方向为 $(\cos2\theta,
\sin2\theta)$；一个法向量为 $(-\sin2\theta,
\cos2\theta)$，故直线为
\[
-\sin2\theta\,(x - \cos\theta) + \cos2\theta\,(y -
\sin\theta) = 0,
\]
而常数项为 $-\sin2\theta\cos\theta +
\cos2\theta\sin\theta = -\sin\theta$：乘以 $-1$，
得 $x\sin2\theta - y\cos2\theta = \sin\theta$。

\textbf{18.} $(a, b, c) = (\sin2\theta, -\cos2\theta,
\sin\theta)$：$a' = 2\cos2\theta$，$b' = 2\sin2\theta$，$c' =
\cos\theta$，$\Delta = 2\sin^22\theta + 2\cos^22\theta = 2$。
由克莱姆法则，再用积化和差公式：
\begin{align*}
x &= \frac{2\sin\theta\sin2\theta +
\cos\theta\cos2\theta}{2}
= \frac{(\cos\theta - \cos3\theta) + \frac12(\cos\theta +
\cos3\theta)}{2} = \frac{3\cos\theta - \cos3\theta}4,\\
y &= \frac{\sin2\theta\cos\theta - 2\cos2\theta\sin\theta}2
= \frac{\frac12(\sin3\theta + \sin\theta) - (\sin3\theta -
\sin\theta)}2 = \frac{3\sin\theta - \sin3\theta}4 :
\end{align*}
此即肾形线，是一条具有两个尖点的闭曲线。

\textbf{19.} 求导并因式分解，其中用到 $\sin3\theta
- \sin\theta = 2\cos2\theta\sin\theta$，$\cos\theta -
\cos3\theta = 2\sin2\theta\sin\theta$：
\[
E'(\theta) = \tfrac34\bigl(\sin3\theta - \sin\theta,\
\cos\theta - \cos3\theta\bigr)
= \tfrac32\sin\theta\,(\cos2\theta, \sin2\theta),
\]
与第 16 题中的反射方向平行：每条
反射光线都与焦散线相切，这正是包络
性质所要求的。$E' = 0$ 恰在 $\sin\theta = 0$ 时成立：
$E(0) = (\tfrac12, 0)$ 与 $E(\pi) = (-\tfrac12, 0)$，即两个
尖点。在直线方程中令 $y = 0$，得 $x\sin2\theta =
\sin\theta$，故 $x = \frac1{2\cos\theta} \to \frac12$，当
$\theta \to 0$ 时，近轴光线交于距圆心 $R/2$ 处 ---
这就是球面镜的焦距。

\textbf{20.} $\norm{E'(\theta)} = \tfrac32\abs{\sin\theta}$，
故弧长为 $\tfrac32\int_0^{2\pi}\abs{\sin\theta}\,
\dd\theta = \tfrac32\cdot4 = 6$。在 $\theta = 0$ 附近，并取
$\cos k\theta = 1 - \tfrac{k^2\theta^2}2 + O(\theta^4)$ 及
$\sin k\theta = k\theta - \tfrac{k^3\theta^3}6 +
O(\theta^5)$:
\[
x - \tfrac12 = \frac{3\theta^2 + O(\theta^4)}{4}
= \tfrac34\theta^2 + O(\theta^4),
\qquad
y = \frac{4\theta^3 + O(\theta^5)}{4}
= \theta^3 + O(\theta^5) :
\]
$p = 2$，$q = 3$，一个沿
坐标轴指向的第一类尖点 --- 即咖啡杯焦散线的亮点。

\textbf{21.} 用
$\Phi(\theta, r) = P_\theta + r\,v_\theta$（沿
每条反射光线的位置）对被照亮的点参数化。雅可比行列式
$\det(\partial_\theta\Phi, \partial_r\Phi) =
\det(P_\theta' + rv_\theta', v_\theta)$ 关于 $r$ 是仿射的，且
恰对唯一的 $r = r_*(\theta)$ 为零 --- 且
$\Phi(\theta, r_*(\theta))$ 即为特征点，因为
此时光线方向与族的变分
线性相关。在包络之外，该映射是局部微分同胚，
而焦散线内侧的点被两条邻近光线
击中（对应两个 $\theta$ 解），外侧的点则不被该族的这一部分中任何光线击中；
在焦散线上两者合并。光
强度与雅可比行列式的绝对值成反比，
因而沿包络趋于无穷：焦散线即
亮曲线，尖点处最亮，此处
退化最为严重。

\textbf{22.} $x' = 1 - \cos t$, $y' = \sin t$, $x'' = \sin
t$，$y'' = \cos t$，故 $x'y'' - y'x'' = \cos t - 1$ 且
$\norm{\gamma'}^2 = 2(1 - \cos t) = 4\sin^2\tfrac t2$;
由\cref{prop:b2:curves:kappaformula}可得，
\[
\kappa(t) = \frac{-(1 - \cos t)}{8\sin^3\tfrac t2}
= \frac{-2\sin^2\tfrac t2}{8\sin^3\tfrac t2}
= -\frac1{4\sin\tfrac t2}
\qquad (0 < t < 2\pi).
\]
取 $T = (\sin\tfrac t2, \cos\tfrac t2)$（将 $\gamma'$
除以 $2\sin\tfrac t2$），并取 $N = (-\cos\tfrac t2, \sin\tfrac
t2)$:
\[
\gamma + \frac1\kappa N
= \gamma - 4\sin\tfrac t2\,\Bigl(-\cos\tfrac t2,\
\sin\tfrac t2\Bigr)
= \bigl(t - \sin t + 2\sin t,\ 1 - \cos t - 2(1 - \cos
t)\bigr),
\]
即 $c(t) = (t + \sin t,\ \cos t - 1)$。代入 $t = u
+ \pi$:
\[
c = \bigl(u + \pi - \sin u,\ -\cos u - 1\bigr)
= \bigl((u - \sin u) + \pi,\ (1 - \cos u) - 2\bigr) :
\]
即摆线 $\gamma(u)$ 经平移 $(\pi, -2)$ 所得。
摆线的渐屈线是一条全等的摆线，
位于其下方一层。

\textbf{23.} $R(t) = 1/\abs{\kappa(t)} = 4\sin\tfrac t2$，故
$R(\pi) = 4$，恰为
\cref{exo:b2:curves:1}中算得的拱弧长 $8$ 的一半。渐屈线的速度 $c'(t) = (1 +
\cos t, -\sin t)$ 在 $t = \pi$ 处为零：尖点为 $c(\pi)
= (\pi, -2)$，位于顶点 $\gamma(\pi) = (\pi,
2)$ 正下方，距离 $4 = R(\pi)$，恰为
该处的密切圆半径的长。

\textbf{24.} 由\cref{exo:b2:curves:6}得 $c'(s) =
-\frac{\kappa'(s)}{\kappa(s)^2}N(s) = R'(s)\,N(s)$，故
$\norm{c'(s)} = \abs{R'(s)}$，且当 $R$ 单调时，
\[
\int_{s_0}^{s_1}\norm{c'(s)}\,\dd s =
\Bigl|\int_{s_0}^{s_1}R'(s)\,\dd s\Bigr| = \abs{R(s_1) -
R(s_0)} .
\]
设 $R$ 递减。沿渐屈线铺设一根绳子，越过
$c(s_0)$ 后再接以从 $c(s_0)$ 到
$\gamma(s_0)$ 的线段（由第 4 题可知，它与渐屈线相切），
剥离至 $c(s)$ 并拉紧时，直线部分
长为 $R(s_0) - \bigl(R(s_0) - R(s)\bigr) = R(s)$
沿法线自 $c(s)$ 指出 --- 恰好落在
$\gamma(s)$ 上：自由端描出原曲线
（“渐伸线”）。惠更斯将一摆悬挂于两块
摆线形挡板之间：绳索缠绕于渐屈线上，故摆锤
描出一条摆线 --- 即等时曲线，其振动周期
与振幅无关。

\textbf{25.} 抛物线的切线：$\Delta = 1$，且 $E' =
(1, t) \neq 0$ 处处成立 --- 光滑包络（抛物线
本身）。滑动梯子：$\Delta = -1$，但 $E' =
\tfrac32\sin 2t\,(-\cos t, \sin t)$ 在象限
端点处为零 --- 即星形线的四个尖点。法线（
抛物线与摆线）：$\Delta = \kappa \neq 0$，且
$E' = R'N$ 恰在曲率取极值处为零
--- 为渐屈线在 $(0,1)$ 与 $(\pi, -2)$ 处的尖点。焦散线：
$\Delta = 2$，且 $E' = \tfrac32\sin\theta\,(\cos2\theta,
\sin2\theta)$ 在 $\theta = 0, \pi$ 处为零 --- 即肾形线的两个尖点，
亦即镜面的焦点。曲率的
极值\emph{必然}在法线族的包络上产生尖点，
因为渐屈线的速度为 $R'N$，这正是
椭圆的渐屈线有四个尖点（即四个顶点）的缘故，
而退化的族（线束、平行曲线族）则属于
机器只输出一个点或完全不输出的情形。
\end{solution}
